\documentclass[12pt]{article}
\title{ECE 11L Lab Report 1}

\author{Lawrence Liu}
\usepackage{graphicx}
\usepackage{amsmath}
\usepackage{subcaption}
\usepackage{array}
\usepackage{multicol}
\begin{document}
\begin{titlepage}
   \begin{center}
        \vspace*{3cm}
        \Huge
        \textbf{EE 3 Final Report}
                
        \vspace{4cm}
        \Large
        \textbf{Lawrence Liu} \\
        \normalsize   
        \textit{lawrencerliu@ucla.edu}\\
        \vspace{1cm}
        \Large
        \textbf{Inesh Chakrabarti} \\
        \normalsize   
        \textit{inesh33@g.ucla.edu}\\

        \vspace{2cm}
            
        \includegraphics[width=0.4\textwidth]{UCLA.png}\\
        \vspace{1cm}
        \large
        Electrical and Computer Engineering\\
        University of California, Los Angeles\\
        United States of America\\
        June 4th 2022
   \end{center}
\end{titlepage}
\section*{Acknowledgements}
We would like to aknowledge and extend my sincere thanks to Dr. Dennis Briggs who made this project possible.
His knoweldge of the subject, guidance, and support through the whole quarter carried us through this project.
We would also like the aknowledge our TAs Dhruv Srinivas and Xin Li for your brilliant comments and suggestions. Furthermore we 
would like to thank from the bottom of our hearts our wonderful mentors Alexis Aleksandrovich Samoylov and Vaibhav Gupta for their guidance and support.
\\\\
Finally we would also want to thank our families, without them none of this would be possible.
\pagebreak
\begin{multicols}{2}
\section*{Introduction}
The goal of this project was to design a closed loop controll system for a small "car" robot to follow
a line. The car was a Texas Instrument (TI) Robotics Systems Learning Kit (RSLK).\\
This car was controlled by a TI MSP-EXP432P401R launchpad microcontroller. This microcontroller is part of TI's MSP432P401x family of 
ultra low power microcontrollers.\\\\ The CPU is a ARM 32-bit Cortex-M4 RISC engine with a frequency of up to 48 MHz. Furthermore the microcontroller
has 256KB of Flash Main Memory, 16KB of Flash Information Memory, and 64KB of SRAM.
\section*{Testing Methodology}
Our development followed two routes, we tried both a basic PID developmental route and a attempt to develop a ML model to controll
the car through Deep Q reinforcement learning.
\\\\
Unfortunately, the Deep Q reinforcement learning could not be made to work. Our code had memory leaks and 
it was difficult working around the microcontroller's small memory size.\\\\
So we result in using a basic PID controller. We realized that since the car is always moving, there is no steady state.
Therefore we did not need the Integral part of the PID controller since there would be no steady state error to eliminate.
Therefore the closed loop transfer function from the input sensor fusion value to the car movement with the car's plant transfer
function being $G(s)$ was
$$\frac{(k_p+k_d s)G(s)}{1+(k_p+k_d s)G(s)}$$
Where $k_p$ and $k_d$ are the proportional and derivative gains respectively. Let us call the output from the controller to be 
$d$, and assume that the car was set to travel at a base speed of
$V_{base}$, then because of our controller, the left and right velocites of the wheels would become
$$V_{left}=V_{base}-d$$
$$V_{right}=V_{base}+d$$
Therefore the parameters we would want modify would be $k_p$ and $k_d$ and the base speed. To determine
the optimal values for these parameters, we would measure whether the car completed the track, and if it did,
the time it took the car to complete the track.
\section*{Analysis}
To facilitate the analysis of the system we developed two metrics to quantify how
the car performed. The first metric was Oscilations, quatified from 5 (full oscilations) to 0 (no oscilations).
. The second was acceptiblity quantify from 10 (fully acceptiblity) to 0 (not acceptable). This criterion was
based on both the speed of the car, the time it took to complete the track, and how close it came to 
completing the track if it did not. \\\\
To analyze how the $k_p$ and $k_d$ values affected these two metrics we ploted a scatter graph of the oscilation and acceptablity
for all the values of $k_p$ and $k_d$ we experimented with, regradless of the base speed.
\begin{center}
\includegraphics*[scale=0.3]{KpKiAll.png}
\end{center}
However this does not take into account our development path. 
We first attempted to controll the car with PID at 50 speed. Once this suceeded we increased the speed to 
100 speed, however the car was unable to keep on the track at this speed. So we focused on making a peicewise PID controller with different base speeds and 
$k_p$ and $k_d$ values.\\\\
 Since we already had a workign set of values for $k_p$ and $k_d$ for 50 speed, we focused on developing another set of
them for 100 speed that could complete everything but the initial chicane. \\\\
Therefore we would want to take into account these factors, such as the velocity, the battery voltage, and whether we were testing over the 
entire track. However this would make the data we would want to visulaize to have more than 2 dimensions, so we would
have to use a dimensionality reduction algorith to reduce the data to 2 dimensions that is plottable. We first try to use
Principal Component Analysis (PCA) to reduce the data to 2 dimensions.
\begin{center}
    \includegraphics*[scale=0.3]{KpKiAllPCA.png}
\end{center}
As we can see, the PCA algorithm did not work well. Therefore lets try to use t-SNE to reduce the data to 2 dimensions.
\begin{center}
    \includegraphics*[scale=0.3]{KpKiAllTSNE.png}
\end{center}
This algorithm seems to work well, and it seems like the data is organized into 2-3 clusters. Therefore let us cluster 
it with K-means at see the average acceptablity for each cluster. K means for 3 clusters classifies the data as the following:
\begin{center}
    \includegraphics*[scale=0.3]{KpKiAllTSNEKMeans.png}
\end{center}
Therefore, K means does sucessfully classify the data into the 3 clusters we would expect
\end{multicols}
\end{document}
